% =====================================================================
\needspace{0.62\textheight}
\appendix
\renewcommand{\thesection}{\Alph{section}}
\section{Formulario del Tema 1}
% =====================================================================
\begingroup
\small
\renewcommand{\arraystretch}{1.85}
\newcommand{\grupo}[1]{\multicolumn{3}{l}{\color{acento}\espaciado[200]{#1}}\\[-2pt]}
\noindent
\begin{tabularx}{\textwidth}{@{}l>{\centering\arraybackslash}Xr@{}}
\toprule
Magnitud o resultado & Expresión & Referencia\\
\midrule
\grupo{Sistemas y estado}
Moles y partículas & $n=m/M,\quad N=nN_A$ & sección~1.2\\
Específica y molar & $x=X/m,\quad\bar X=X/n$ & ec.~\eqref{eq:molar}\\
Temperatura absoluta & $T/\mathrm K=t/{}^\circ\mathrm C+273{,}15,\quad\Delta T=\Delta t$ & sección~3\\
Temperatura empírica & $\theta_1(p_1,V_1)=\theta_2(p_2,V_2)=\theta_3(p_3,V_3)$ & ec.~\eqref{eq:theta}\\
\midrule
\grupo{Ecuaciones de estado y coeficientes}
Gas ideal & $pV=nRT,\quad\alpha=1/T,\quad\kappa_T=1/p,\quad\beta=1/T$ & ec.~\eqref{eq:gasideal}\\
Coeficientes & $\alpha=\frac1V\pdi{V}{T}{p},\ \ \kappa_T=-\frac1V\pdi{V}{p}{T},\ \ \beta=\frac1p\pdi{p}{T}{V}$ & ec.~\eqref{eq:coef}\\
Relación $\alpha$, $\beta$, $\kappa_T$ & $\alpha=p\beta\kappa_T$ & ec.~\eqref{eq:apbk}\\
Recíproca y cíclica & $\pdi{x}{y}{z}=1/\pdi{y}{x}{z},\quad\pdi{x}{y}{z}\pdi{y}{z}{x}\pdi{z}{x}{y}=-1$ & ec.~\eqref{eq:reciclica}\\
Exactitud & $M\,\dd x+N\,\dd y$ exacta $\iff\partial_yM=\partial_xN$ & ec.~\eqref{eq:exacta}\\
Estabilidad mecánica & $\kappa_T>0\iff\pdi{p}{V}{T}<0$ & obs.~\ref{obs:kneg}\\
\midrule
\grupo{Ecuación diferencial de estado}
Diferenciales & $\dd\ln V=\alpha\,\dd T-\kappa_T\,\dd p,\quad\dd p=p\beta\,\dd T-\dfrac{\dd V}{\kappa_TV}$ & ec.~\eqref{eq:resumen}\\
Regla del cociente & $(f/g)'=(f'g-fg')/g^2$ & ec.~\eqref{eq:cociente}\\
Logaritmos & $a\ln x=\ln x^a,\quad\ln x-\ln y=\ln(x/y),\quad\e^{\ln X+y}=X\e^y$ & ec.~\eqref{eq:logs}\\
Coeficientes constantes & $\ln(V_2/V_1)=\alpha\Delta T-\kappa_T\Delta p$ & ec.~\eqref{eq:const}\\
Volumen constante & $\pdi{p}{T}{V}=\alpha/\kappa_T=p\beta$ & probl.~1.7\\
Dilatación lineal & $\alpha=3\alpha_l,\quad\Delta L\simeq\alpha_lL_0\Delta T$ & probl.~1.5, 1.6\\
Politrópico & $pV^n=\cte,\quad n=\ln(p_2/p_1)/\ln(V_1/V_2)$ & probl.~1.3, 1.4\\
Punto crítico (Dieterici) & $T_c=a/4Rb,\quad\Vm_c=2b,\quad p_c=a/(4\e^2b^2)$ & probl.~1.14\\
\bottomrule
\end{tabularx}
\endgroup

% =====================================================================
\begingroup
\titleformat{\section}{\Large\bfseries\color{navy}}{}{0em}{}[\vspace{2pt}{\color{navy}\titlerule[0.9pt]}]
\renewcommand{\refname}{Referencias}
\makeatletter
\renewenvironment{thebibliography}[1]
  {\section*{\refname}\addcontentsline{toc}{section}{\refname}\markboth{\refname}{}%
   \list{\@biblabel{\@arabic\c@enumiv}}{\settowidth\labelwidth{\@biblabel{#1}}\leftmargin\labelwidth
     \advance\leftmargin\labelsep\itemsep 4pt\usecounter{enumiv}\let\p@enumiv\@empty
     \renewcommand\theenumiv{\@arabic\c@enumiv}}\sloppy}
  {\def\@noitemerr{\@latex@warning{Empty `thebibliography' environment}}\endlist}
\makeatother
\clearpage
\begin{thebibliography}{9}
\bibitem{gonzalez1} J.~González Ayala. \emph{Termodinámica I}, capítulo 1: Conceptos básicos, y Lista de
  problemas \#1. Apuntes de la asignatura. Grado en Física, Universidad de Salamanca, 2026.
\bibitem{gonzalez2} J.~González Ayala. \emph{Tema 1. Conceptos básicos: de la descripción macroscópica a
  la ecuación de estado}. Transparencias. Universidad de Salamanca, curso 2026--27.
\bibitem{parcial} \emph{Termodinámica I. Prueba de evaluación continua (temas 1--3)}. Universidad de
  Salamanca, 31 de octubre de 2025.
\bibitem{callen} H.~B.~Callen. \emph{Thermodynamics and an Introduction to Thermostatistics}, 2.ª ed. Wiley, 1985.
\bibitem{zemansky} M.~W.~Zemansky y R.~H.~Dittman. \emph{Calor y termodinámica}, 6.ª ed. McGraw-Hill, 1990.
\bibitem{aguilar} J.~Aguilar Peris. \emph{Curso de termodinámica}. Alhambra, 1981.
\end{thebibliography}
\endgroup
